\chapter{Policy-Gradient Foundations for LLM Reinforcement Learning}\label{ch:pg}

Every reinforcement-learning method in the rest of this part (PPO with a learned critic in
\cref{ch:rlhf}, GRPO and its descendants in \cref{ch:grpo,ch:grpo2}, and the
distributed systems that run them in \cref{ch:rlsys}) is built from the same small set
of ideas: the score-function (REINFORCE) gradient, baselines that reduce its variance,
normalisation choices that quietly change what is being optimised, and importance weights
that let us reuse samples drawn from an older policy. Preference methods that avoid RL
altogether (DPO, \cref{ch:dpo}) are easiest to understand as a deliberate step away from
these ideas.

This chapter builds those foundations slowly and carefully, with two audiences in mind.
For the theory, we derive each estimator, state exactly what it is unbiased for, and
measure its variance on problems small enough to compute exactly. For the systems, we
spend as much time on \emph{how the batch is laid out in memory}, \emph{how log-probabilities
are computed without materialising a 60\,GB tensor}, and \emph{how to make micro-batched
gradients bit-for-bit equal to the full-batch gradient} as on the mathematics. Both kinds of
mistake are silent: an RL run with a wrong normaliser or a misaligned response mask still
trains, it just optimises something else. The chapter ends with the engineering loop that
wires a rollout engine, a reward function, and a trainer into a working system. That loop
is exactly what the CS336 reasoning-RL assignment (\cref{ch:proj5}) asks you to build at
scale.

\begin{figure}[t]
\centering
\resizebox{\linewidth}{!}{%
\begin{tikzpicture}[st/.style={draw,rounded corners=3pt,minimum width=2.6cm,minimum height=1.0cm,
  align=center,font=\footnotesize,text width=2.45cm},>={Stealth[length=2mm]}]
\node[st,fill=fdata!15] (s1) at (0,0) {\textbf{1} RL in LM\\notation};
\node[st,fill=fdata!15] (s2) at (2.95,0) {\textbf{2} REINFORCE\\derivation};
\node[st,fill=fcompute!15] (s3) at (5.9,0) {\textbf{3} Baselines \&\\variance};
\node[st,fill=fcompute!15] (s4) at (8.85,0) {\textbf{4} Normalisation\\\& reweighting};
\node[st,fill=fmemory!20] (s5) at (0,-2.0) {\textbf{5} Batch layout,\\log-probs, entropy};
\node[st,fill=fmemory!20] (s6) at (2.95,-2.0) {\textbf{6} Aggregation \&\\accumulation};
\node[st,fill=fcomm!15] (s7) at (5.9,-2.0) {\textbf{7} Off-policy:\\IS, TRPO, PPO};
\node[st,fill=fcomm!15] (s8) at (8.85,-2.0) {\textbf{8} RFT / expert\\iteration};
\node[st,fill=labpurple!15,minimum width=11.4cm,text width=11.2cm] (s9) at (4.425,-3.6)
  {\textbf{9} The engineering loop: rollout engine $\to$ reward $\to$ trainer $\to$ weight sync, with logging};
\foreach \a/\b in {s1/s2,s2/s3,s3/s4,s5/s6,s6/s7,s7/s8} \draw[->,thick] (\a)--(\b);
\draw[->,thick] (s4.south) -- ++(0,-0.25) -| (s5.north);
\draw[->,thick] (s8.south) -- (s8.south |- s9.north);
\end{tikzpicture}}
\caption{Roadmap. Stages 1--4 are the mathematics of the estimator; 5--6 make it correct on
real hardware; 7--8 cover reusing samples and the simplest alternative to RL; 9 assembles
everything into the system that later chapters scale up.}
\label{fig:pg:roadmap}
\end{figure}

% =====================================================================
\section{Reinforcement learning, and why language models fit it}\label{sec:pg:rl}
\covers{PO §3; PO §4.1; CS336 L16}

\paragraph{The classical setting.} A reinforcement-learning problem is a Markov decision
process (MDP): a set of states $\mathcal{S}$, a set of actions $\mathcal{A}$, an initial-state
distribution $\rho(s_0)$, a transition kernel $P(s_{t+1}\mid s_t,a_t)$, a reward function
$r(s_t,a_t)$, and a discount $\gamma\in[0,1]$. An \emph{agent} with a stochastic
\emph{policy} $\pitheta(a\mid s)$ produces a \emph{trajectory}
$\tau=(s_0,a_0,r_0,s_1,a_1,r_1,\dots,s_T)$ by alternating actions and transitions. Its goal is
to maximise the expected \emph{return}
\begin{equation}
J(\theta)=\E_{\tau\sim\pitheta}\big[R(\tau)\big],\qquad R(\tau)=\sum_{t=0}^{T-1}\gamma^t r_t .
\label{eq:pg:J}
\end{equation}
Three derived quantities appear everywhere. The \emph{state value}
$V^\pi(s)=\E[\sum_{k\ge0}\gamma^k r_{t+k}\mid s_t=s]$ is the expected return from a state. The
\emph{action value} $Q^\pi(s,a)$ is the same quantity after first taking action $a$. The
\emph{advantage} $A^\pi(s,a)=Q^\pi(s,a)-V^\pi(s)$ says how much better $a$ is than the
policy's average behaviour in $s$.

\paragraph{Three families of algorithms.} \emph{Value-based} methods (Q-learning, DQN) learn
$Q$ and act greedily with respect to it; the policy is implicit. \emph{Policy-based} methods
(REINFORCE) parameterise $\pitheta$ directly and follow $\nabla_\theta J$. \emph{Actor--critic}
methods (A2C, TRPO, PPO) do both: an actor (the policy) is updated with advantages computed
from a critic (a learned $V$). For language models the action space is the vocabulary,
$|\mathcal{A}|\approx10^5$, and the state is an entire text prefix, so learning $Q(s,a)$
for every token is hopeless; every successful LLM-RL method is policy-based or
actor--critic. Value-based ideas reappear only as critics or as process reward models.

\begin{figure}[t]
\centering
\resizebox{\linewidth}{!}{%
\begin{tikzpicture}[b/.style={draw,rounded corners=2pt,align=center,font=\small,minimum height=9mm},
  >={Stealth[length=2mm]}]
\begin{scope}
\node[font=\small\bfseries] at (2.2,2.2) {(a) classical RL};
\node[b,fill=fcompute!15,text width=19mm] (ag) at (0,0) {agent\\$\pitheta(a\mid s)$};
\node[b,fill=fdata!15,text width=21mm] (env) at (4.4,0) {environment\\$P(s'\mid s,a)$, $r$};
\draw[->,thick] (ag.north) to[out=50,in=130] node[above,font=\scriptsize]{action $a_t$} (env.north);
\draw[->,thick] (env.south) to[out=-130,in=-50] node[below,font=\scriptsize]{state $s_{t+1}$, reward $r_t$} (ag.south);
\end{scope}
\begin{scope}[xshift=8.3cm]
\node[font=\small\bfseries] at (2.6,2.2) {(b) a language model generating a response};
\node[b,fill=fcompute!15,text width=26mm] (lm) at (0,0) {LM\\$\pitheta(y_t\mid x,y_{<t})$};
\node[b,fill=white,text width=24mm] (ap) at (5.0,0) {append token\\{\scriptsize (deterministic)}};
\node[b,fill=napkin!15,text width=20mm] (ver) at (5.0,-2.1) {verifier /\\reward model};
\draw[->,thick] (lm.north) to[out=50,in=130] node[above,font=\scriptsize]{token $y_t$} (ap.north);
\draw[->,thick] (ap.south west) to[out=-150,in=-30] node[below,font=\scriptsize,pos=0.5]{prefix $(x,y_{\le t})$} (lm.south east);
\draw[->,thick,dashed,napkin] (ap.south)--node[right,font=\scriptsize]{at EOS: whole $y$}(ver.north);
\draw[->,thick,dashed,napkin] (ver.west) -| node[pos=0.3,below,font=\scriptsize]{one scalar $r(x,y)$} (lm.south);
\end{scope}
\end{tikzpicture}}
\caption{(a) The textbook agent--environment loop. (b) The same loop for an LLM: the
``environment'' is just concatenation, so transitions are deterministic and known, and the
only external signal is a single reward at the end of the response. Everything that makes
LLM-RL different from robotics or games follows from (b): enormous action spaces, perfect
simulators for free, and sparse terminal rewards.}
\label{fig:pg:loops}
\end{figure}

\subsection{The notation map between classical RL and language modelling}\label{sec:pg:notation}
\covers{PO §5 (notation map); CS336 A5 §2.3, §4.1.1--4.1.3}

The two communities use different words for the same objects, and papers mix them freely.
\Cref{tab:pg:notation} is the dictionary we use for the rest of the book; it extends the
short glossary of \cref{tab:grpo:vocab}.

{\small
\begin{longtable}{@{}p{0.22\linewidth}p{0.25\linewidth}p{0.47\linewidth}@{}}
\caption{Classical RL objects and their language-model counterparts.}\label{tab:pg:notation}\\
\toprule
\textbf{Classical RL} & \textbf{Language modelling} & \textbf{Notes for LLMs}\\
\midrule\endfirsthead
\toprule
\textbf{Classical RL} & \textbf{Language modelling} & \textbf{Notes for LLMs}\\
\midrule\endhead
initial-state distribution $\rho(s_0)$ & prompt distribution / dataset & $x\sim\rho$; often a fixed set of problems, sampled without replacement per epoch\\
state $s_t$ & prompt plus prefix $(x,y_{<t})$ & the entire text so far; never summarised\\
action $a_t$ & next token $y_t$ & $|\mathcal{A}|=V\approx 3\times10^4$--$2.5\times10^5$\\
policy $\pi(a\mid s)$ & the LM's next-token distribution & $\softmax(f_\theta(x,y_{<t})/\tau)$; temperature, top-$p$ and top-$k$ change the \emph{sampling} policy\\
transition $P(s'\mid s,a)$ & append the token: $s_{t+1}=(s_t,a_t)$ & deterministic; tool calls in agentic RL make it stochastic again (\cref{ch:rlsys})\\
terminal state, horizon $T$ & EOS, stop string, or token budget & horizon is variable, up to the max generation length $L$\\
trajectory, episode, rollout & response, completion, sample $y$ & one prompt, one sampled answer\\
reward $r(s_t,a_t)$ & usually zero except at the last token & verifier (0/1), reward model (real), or shaped per token (KL penalty, \cref{ch:rlhf})\\
return $R(\tau)$ & $r(x,y)$ & with terminal reward and $\gamma=1$, every token sees the same return\\
discount $\gamma$ & usually $1$ & episodes are finite and short in decision count compared with games\\
value $V(s_t)$ & expected final reward from a prefix & learned by a critic (PPO) or replaced by group statistics (GRPO)\\
action value $Q(s_t,a_t)$ & expected reward after committing to $y_t$ & almost never learned explicitly for LLMs\\
advantage $A_t$ & how much better $y_t$ (or $y$) was than average & per sequence (GRPO, RLOO) or per token (PPO with GAE)\\
behaviour policy $\beta$ / $\piold$ & the weights the rollout engine used & differs from $\pitheta$ after the first update, and numerically even before (\cref{ch:rlsys})\\
reference policy $\piref$ & the SFT or base model & anchor for KL regularisation; frozen\\
critic & value model (an LM with a scalar head) & PPO only; a second network as big as the policy\\
reward model $r_\phi$ & preference-trained scorer & RLHF; replaced by a program in RLVR\\
batch of episodes & $B$ prompts $\times$ $G$ samples & $G$ is the ``group size''\\
contextual bandit view & one action $=$ one whole response & valid when reward is terminal and $\gamma=1$; GRPO and DPO effectively use it\\
\bottomrule
\end{longtable}}

\paragraph{Token MDP or contextual bandit?} \Cref{tab:pg:notation} hides a modelling choice.
With a terminal reward and no discount, we can treat the whole response as \emph{one} action
taken in context $x$: a contextual bandit whose action space is all strings. Or we can treat
each token as an action in a token-level MDP. The policy gradient is identical in both views,
because $\log\pitheta(y\mid x)=\sum_t\log\pitheta(y_t\mid x,y_{<t})$. They differ in
\emph{credit assignment}: a token-level critic can give different advantages to different
tokens of the same response, while the bandit view gives every token the same advantage.
\Cref{fig:pg:tree} shows both readings of one rollout.

\begin{figure}[t]
\centering
\begin{tikzpicture}[tok/.style={draw,rounded corners=1.5pt,minimum width=10mm,minimum height=6mm,font=\scriptsize},
  >={Stealth[length=1.6mm]}]
\node[tok,fill=fdata!20] (x) at (0,0) {$x$};
\node[tok,fill=fcompute!15] (a1) at (1.6,0) {\texttt{12}};
\node[tok,fill=fcompute!15] (a2) at (3.2,0) {\texttt{+}};
\node[tok,fill=fcompute!15] (a3) at (4.8,0) {\texttt{30}};
\node[tok,fill=fcompute!15] (a4) at (6.4,0) {\texttt{=42}};
\node[tok,fill=fcompute!15] (a5) at (8.0,0) {\texttt{EOS}};
\node[tok,fill=napkin!25] (r) at (9.7,0) {$r{=}1$};
\node[tok,fill=black!5] (b1) at (1.6,-1.0) {\texttt{The}};
\node[tok,fill=black!5] (b2) at (3.2,-1.0) {\texttt{*}};
\node[tok,fill=black!5] (b3) at (4.8,-1.0) {\texttt{3}};
\draw[->] (x)--(a1); \draw[->] (a1)--(a2); \draw[->] (a2)--(a3); \draw[->] (a3)--(a4);
\draw[->] (a4)--(a5); \draw[->,dashed] (a5)--(r);
\draw[->,black!50] (x)--(b1); \draw[->,black!50] (a1)--(b2); \draw[->,black!50] (a2)--(b3);
\node[font=\scriptsize,text=black!60,anchor=west] at (5.5,-1.0) {tokens the sampler did not pick};
\draw[decorate,decoration={brace,amplitude=4pt},thick,keygreen] (1.05,0.45)--(8.55,0.45)
  node[midway,above=5pt,font=\scriptsize,keygreen] {bandit view: one action $y=$\,``\texttt{12+30=42}'' with reward $r(x,y)$};
\draw[decorate,decoration={brace,amplitude=4pt,mirror},thick,sysblue] (1.05,-1.5)--(8.55,-1.5)
  node[midway,below=5pt,font=\scriptsize,sysblue] {token MDP: five decisions in five states $(x,y_{<t})$; reward only after EOS};
\end{tikzpicture}
\caption{One rollout read two ways. The sampled path (blue) is a sequence of token decisions
in a deterministic tree; the reward arrives only after EOS. Group-based methods treat the
whole path as a single bandit action; critic-based methods estimate a value at every node.}
\label{fig:pg:tree}
\end{figure}

\subsection{Online and offline, on-policy and off-policy}\label{sec:pg:online}
\covers{PO §3.1}

Two distinctions organise the whole landscape of post-training methods, and they are often
confused.

\emph{Online versus offline} is about \textbf{where the data come from}. An online method
samples fresh responses from the model during training; an offline method learns from a
fixed dataset collected beforehand. SFT and DPO on a static preference set are offline;
PPO-RLHF and GRPO are online.

\emph{On-policy versus off-policy} is about \textbf{whether the estimator assumes the data
came from the current weights}. Writing $\beta$ for the behaviour policy that produced the
data, an off-policy method must correct for $\beta\ne\pitheta$, typically with importance
weights:
\begin{equation}
J(\theta)=\E_{y\sim\pitheta}[r(y)]=\E_{y\sim\beta}\!\Big[\frac{\pitheta(y)}{\beta(y)}\,r(y)\Big].
\label{eq:pg:is}
\end{equation}
Offline data are always off-policy. Online data can still be slightly off-policy: PPO takes
several gradient steps on each batch, asynchronous RL systems train on rollouts produced by
weights that are a few versions old (\cref{ch:rlsys}), and even a ``synchronous'' system
samples with an inference engine whose numerics differ from the trainer's.

\begin{table}[ht]
\centering\small
\begin{tabularx}{\linewidth}{@{}lllX@{}}
\toprule
Method & Data & Policy & What corrects for staleness\\
\midrule
SFT on demonstrations & offline & off & nothing (it is not RL)\\
Rejection fine-tuning (RFT) & online, iterated & on per round & resampling every round\\
REINFORCE / RLOO & online & on & one update per batch\\
GRPO (one step per batch) & online & on & nothing needed\\
PPO / GRPO ($k$ steps per batch) & online & mildly off & clipped token ratios\\
Asynchronous RL & online & off by $1$--$n$ versions & truncated IS, masking (\cref{ch:rlsys})\\
DPO, IPO, KTO & offline (or iterated) & off & a KL-anchored closed form (\cref{ch:dpo})\\
\bottomrule
\end{tabularx}
\caption{Where common LLM post-training methods sit on the two axes.}
\label{tab:pg:onoff}
\end{table}

% =====================================================================
\section{The REINFORCE gradient, from first principles}\label{sec:pg:reinforce}
\covers{PO §3.2; PO §4.2; CS336 A5 §4.1.4}

We want $\nabla_\theta J$ for $J(\theta)=\E_{x\sim\rho}\E_{y\sim\pitheta(\cdot\mid x)}[r(x,y)]$.
The reward is a black box: a regex, a SymPy check, a unit-test runner, or a reward model whose
gradients we do not want. The \emph{score-function} identity moves the gradient off the
reward and onto the log-probability:
\begin{align}
\nabla_\theta \E_{y\sim\pitheta}[r(y)]
 &=\sum_y r(y)\,\nabla_\theta\pitheta(y)
  =\sum_y r(y)\,\pitheta(y)\,\nabla_\theta\log\pitheta(y)\nonumber\\
 &=\E_{y\sim\pitheta}\big[r(y)\,\nabla_\theta\log\pitheta(y)\big].
\label{eq:pg:score}
\end{align}
(We suppress $x$; the outer expectation over prompts commutes with $\nabla_\theta$.) The
only facts used are $\nabla\pi=\pi\nabla\log\pi$ and that the sampling distribution, not
the reward, depends on $\theta$. The sum over $y$ ranges over every string, but we never
evaluate it: \cref{eq:pg:score} is an expectation, so we can estimate it from samples. With
$B$ prompts and $G$ samples each,
\begin{equation}
\hat g=\frac{1}{BG}\sum_{i=1}^{B}\sum_{j=1}^{G} r(x_i,y_{ij})\sum_{t=1}^{|y_{ij}|}
\nabla_\theta\log\pitheta(y_{ij,t}\mid x_i,y_{ij,<t}).
\label{eq:pg:reinforce}
\end{equation}

\paragraph{Reward-to-go.} In a general MDP, an action cannot affect rewards received before
it, so the return in \cref{eq:pg:score} can be replaced, at token $t$, by the reward-to-go
$\sum_{k\ge t}\gamma^{k-t}r_k$ without changing the expectation. This reduces variance for
free. With a single terminal reward, the reward-to-go at every token \emph{is} the final
reward, so for verifier-based RL there is nothing to gain. The trick matters again when
rewards are shaped per token, for example with a per-token KL penalty (\cref{ch:rlhf}).

\begin{keyidea}[the pseudo-loss is not a loss]
Autograd computes gradients of scalars, so implementations build a \emph{surrogate} whose
gradient equals $\hat g$:
\[
\mathcal{L}_{\text{PG}}(\theta)=-\frac{1}{BG}\sum_{i,j}\operatorname{sg}\!\big[A_{ij}\big]
\sum_t\log\pitheta(y_{ij,t}\mid\cdot),
\]
where $\operatorname{sg}$ is stop-gradient and $A_{ij}$ is the reward or advantage. The
minus sign turns gradient ascent into the descent that optimisers perform. The
\emph{value} of $\mathcal{L}_{\text{PG}}$ is meaningless as a progress signal: it depends on
the batch's advantages, can be negative, and may rise while the policy improves. Log rewards,
accuracies and entropies, not this ``loss''.
\end{keyidea}

\Cref{lst:pg:surrogate} shows the whole estimator in the book's API. Everything else in this
chapter changes either the weight multiplying $\log\pi$ or how per-token terms are summed.

\begin{lstlisting}[caption={The REINFORCE surrogate for a batch of fixed-length responses (from \code{llmsys.pg}).},label={lst:pg:surrogate}]
def reinforce_step(model, prompts, responses, rewards, opt):
    """prompts: (N, P) ids, responses: (N, k) ids, rewards: (N,) floats."""
    full = torch.cat([prompts, responses], 1)
    logits = model(full[:, :-1])                       # predict tokens 1..T-1
    P = prompts.shape[1]
    logp, _ = token_logprobs_entropy(logits[:, P - 1:], full[:, P:])   # (N, k)
    loss = -(rewards[:, None].detach() * logp).mean()  # surrogate; value is not meaningful
    opt.zero_grad(set_to_none=True)
    loss.backward()
    opt.step()
\end{lstlisting}

% =====================================================================
\section{Baselines and the variance of the estimator}\label{sec:pg:baselines}
\covers{PO §3.2; CS336 A5 §4.1.5; CS336 L16}

\Cref{eq:pg:reinforce} is unbiased but noisy. Its noise has a simple source. With binary
rewards, an incorrect response contributes nothing and a correct response pushes up
\emph{every} token in it, including the irrelevant ones, by an amount that depends on how
lucky the sample was. A \emph{baseline} $b$ subtracted from the reward fixes part of this.

\begin{proposition}[baselines preserve the expectation]\label{prop:pg:baseline}
If $b$ does not depend on the sampled response $y$ (it may depend on $x$, on $\theta$, or on
other, independent samples), then
$\E_{y\sim\pitheta}\big[(r(y)-b)\nabla_\theta\log\pitheta(y)\big]=\nabla_\theta J$.
\end{proposition}
\emph{Proof.} $\E_y[b\nabla\log\pitheta(y)]=b\sum_y\nabla\pitheta(y)=b\,\nabla\sum_y\pitheta(y)=b\nabla1=0$. \qed

The baseline does change the \emph{variance}. Write $g(y)=\nabla_\theta\log\pitheta(y)$ and
consider the total variance (the trace of the covariance) of the one-sample estimator
$(r-b)g$:
\begin{equation}
\operatorname{Var}_b=\E\big[(r-b)^2\|g\|^2\big]-\|\nabla J\|^2 .
\end{equation}
The second term does not depend on $b$, and the first is a quadratic in $b$. Setting its
derivative to zero gives the variance-optimal constant baseline
\begin{equation}
b^\star=\frac{\E\big[r\,\|g\|^2\big]}{\E\big[\|g\|^2\big]},
\label{eq:pg:bstar}
\end{equation}
a reward average \emph{weighted by the squared norm of the score}. It is not the mean reward,
although it is often close to it. Outcomes whose score vectors are large, typically
low-probability actions, pull $b^\star$ towards their rewards.

\paragraph{A worked example you can check exactly.} Take a three-armed bandit with
$\pi=\softmax(0.5,\,0,\,-0.5)=(0.506,\,0.307,\,0.186)$ and deterministic rewards
$r=(1.0,\,0.2,\,0.0)$. The score of arm $a$ with respect to the logits is $e_a-\pi$, so every
term in \cref{eq:pg:bstar} is a three-term sum. The mean reward is $\E[r]=0.568$, but the
optimal baseline is $b^\star=0.384$. The arm most likely to be sampled has the smallest score
norm, so it counts least in \cref{eq:pg:bstar}. \Cref{fig:pg:variance} plots the exact
variance as a function of $b$. No baseline gives $0.126$, the mean-reward baseline gives
$0.056$, and $b^\star$ gives $0.036$, a $3.5\times$ reduction for free. A baseline of $1.0$
(``everything is a failure'') is \emph{worse} than none. Lab~24 reproduces every number by
sampling.

\begin{figure}[t]
\centering
\begin{tikzpicture}
\begin{axis}[width=0.78\linewidth,height=5.2cm,xlabel={constant baseline $b$},
  ylabel={total variance of $\hat g$ (1 sample)},xmin=-0.2,xmax=1.2,ymin=0,ymax=0.46,
  grid=major,grid style={black!10},tick label style={font=\scriptsize},label style={font=\small}]
\addplot[thick,sysblue] coordinates {(-0.20,0.2454) (-0.15,0.2110) (-0.10,0.1797) (-0.05,0.1515)
 (0.00,0.1264) (0.05,0.1043) (0.10,0.0853) (0.15,0.0693) (0.20,0.0565) (0.25,0.0467) (0.30,0.0400)
 (0.35,0.0363) (0.40,0.0358) (0.45,0.0383) (0.50,0.0438) (0.55,0.0525) (0.60,0.0642) (0.65,0.0790)
 (0.70,0.0968) (0.75,0.1178) (0.80,0.1418) (0.85,0.1688) (0.90,0.1990) (0.95,0.2322) (1.00,0.2685)
 (1.05,0.3079) (1.10,0.3503) (1.15,0.3958) (1.20,0.4444)};
\addplot[only marks,mark=*,warnred] coordinates {(0.384,0.0356)};
\addplot[only marks,mark=square*,keygreen] coordinates {(0.568,0.0563)};
\addplot[only marks,mark=triangle*,napkin] coordinates {(0,0.1264)};
\node[font=\scriptsize,warnred,anchor=north] at (axis cs:0.384,0.03) {$b^\star{=}0.384$};
\node[font=\scriptsize,keygreen,anchor=south west] at (axis cs:0.58,0.058) {$b{=}\E[r]{=}0.568$};
\node[font=\scriptsize,napkin,anchor=south west] at (axis cs:0.01,0.13) {no baseline};
\end{axis}
\end{tikzpicture}
\caption{Exact variance of the one-sample policy gradient for the three-armed bandit of the
text, as a function of a constant baseline. The curve is a parabola with its minimum at the
score-weighted mean reward $b^\star$, not at the mean reward. Computed exactly with
\code{llmsys.pg.bandit\_grad\_moments}; Lab~24 confirms it by Monte Carlo.}
\label{fig:pg:variance}
\end{figure}

\subsection{Baselines you can actually compute for an LLM}\label{sec:pg:practical}

For a language model, $\|g\|^2$ is a norm over billions of parameters, so \cref{eq:pg:bstar}
is a guide, not a recipe. Practical baselines estimate the expected reward \emph{of the
prompt}, $\eta(x)=\E_{y\sim\pitheta(\cdot\mid x)}r(x,y)$, which differs wildly across prompts:
one problem is solved 95\% of the time, the next 3\%. A single global baseline cannot track
that, so all practical choices are per prompt.

\begin{description}
\item[Learned value baseline.] A critic $V_\psi(x)$ (or $V_\psi(x,y_{<t})$ per token) trained
  by regression on observed returns. This is PPO's choice, which \cref{ch:rlhf} develops with
  GAE. It costs a second network the size of the policy.
\item[Group mean.] Sample $G$ responses per prompt and subtract their mean $\bar r$. The
  sample includes $y$ itself, so the baseline is \emph{not} independent of $y$, and the
  expectation picks up a factor $(G-1)/G$: a harmless rescaling of the step size, not a
  change of direction.
\item[Leave-one-out (RLOO).] Subtract the mean of the \emph{other} $G-1$ samples,
  $b_j=\frac{1}{G-1}\sum_{k\ne j}r_k$. Now $b_j$ is independent of $y_j$, so
  \cref{prop:pg:baseline} applies exactly and the estimator is unbiased (Kool et al., 2019;
  Ahmadian et al., 2024). Algebraically, $r_j-b_j=\frac{G}{G-1}(r_j-\bar r)$: RLOO \emph{is}
  the group-mean estimator, rescaled.
\end{description}
The rescaling identity is worth internalising: two papers that ``compare'' the group-mean
and leave-one-out baselines under Adam are comparing the same algorithm, because Adam is
invariant to a constant rescaling of the gradient. The real differences between GRPO
variants lie elsewhere: in dividing by statistics that \emph{do} depend on the group, and in
how tokens are weighted. These are the subject of the next section.

\begin{systemslens}[groups make rollouts a structured workload]
A group baseline forces the rollout engine to sample $G$ continuations of the \emph{same}
prompt. That is good news for the systems: the prompt's KV cache can be computed once and
shared copy-on-write by all $G$ samples (\cref{ch:scaling,ch:serving}), so prefill cost is
paid $B$ times rather than $BG$ times. It is bad news for load balance. All $G$ samples must
finish before the group's advantages exist, so one long sample holds back the other
$G-1$ in the trainer's queue. This is the long-tail problem of \cref{ch:rlsys}, and the reason
production systems schedule rollouts per group, not per sequence.
\end{systemslens}

\begin{napkinmath}[how many groups are wasted?]
With binary rewards and per-prompt success probability $p$, a group of $G$ samples has
\emph{zero} variance, and therefore zero advantage for every member, with probability
$p^G+(1-p)^G$. For $G=8$, that is $0.66$ at $p=0.05$, $0.43$ at $p=0.1$ or $p=0.9$, $0.058$
at $p=0.3$, and $0.008$ at $p=0.5$. On a dataset that is mostly too easy or too hard, more than
half of the generated tokens can carry exactly zero gradient. Two consequences follow. For
the systems: skip zero-advantage sequences in the training forward and backward passes, but
\emph{keep the normaliser fixed} so the gradient is unchanged (\cref{sec:pg:aggregation}).
For the algorithm: filter or resample prompts to keep $p$ near the middle (dynamic sampling,
\cref{ch:grpo2}).
\end{napkinmath}
% =====================================================================
\section{Normalisation changes the objective}\label{sec:pg:normalisation}
\covers{CS336 A5 §4.1.6--4.1.7, §5; CS336 L16}

Baselines are the only modification that provably leaves the expected gradient alone. GRPO
(\cref{ch:grpo}) makes two more: it divides each group's centred rewards by the group's
standard deviation, and it averages each response's token terms over the response's length.
Both are usually described as stability tricks. Both also change \emph{which objective the
estimator ascends}, and it pays to see exactly how.

\subsection{Scaling advantages reweights prompts}\label{sec:pg:reweight}

Any per-group factor that depends on the group's rewards acts, in the limit of large groups,
as a \emph{weight on prompts}. Here is the general argument, worked through for a scheme that
no standard algorithm uses, so you can apply the method to the ones that do.

Suppose we form advantages $A_{ij}=(r_{ij}-\bar r_i)\,\kappa(\bar r_i)$, where
$\bar r_i$ is the group mean and $\kappa$ is any smooth function. Take
$\kappa(\mu)=1-\mu$ as a concrete example: a scheme that shrinks updates on prompts the model
already solves. As $G\to\infty$ the group mean converges to the prompt's success rate,
$\bar r_i\to\eta(x_i)=\E_{y\sim\pitheta}r(x_i,y)$, so for one prompt
\begin{equation}
\frac1G\sum_j (r_{ij}-\bar r_i)\,\kappa(\bar r_i)\,\nabla\log\pitheta(y_{ij})
\;\longrightarrow\;\kappa(\eta(x_i))\,\E_y\!\big[(r-\eta(x_i))\nabla\log\pitheta(y)\big]
=\kappa(\eta(x_i))\,\nabla_\theta\eta(x_i).
\end{equation}
The last step is \cref{prop:pg:baseline}: in the limit the baseline no longer depends on any
single sample. Averaging over prompts, the estimator ascends
\begin{equation}
J_w(\theta)=\E_{x\sim\rho}\big[\,w(x)\,\eta_\theta(x)\big],\qquad
w(x)=\operatorname{sg}\big[\kappa(\eta(x))\big],
\end{equation}
a \emph{difficulty-reweighted} success rate in which the weights are treated as constants
when differentiating. For our example, $w(x)=1-\eta(x)$: a prompt solved 90\% of the time
counts one ninth as much as a prompt solved 10\% of the time. Nothing about this objective is
wrong; it is simply a different objective from accuracy, and it was chosen implicitly.

The same lens explains the published variants. Dividing by the group standard deviation
(GRPO) makes the weight depend on the Bernoulli spread $\sqrt{\eta(1-\eta)}$, and critics of
that choice (Liu et al., 2025, ``Dr.\ GRPO'') argue that it gives too much weight to prompts
that are almost always solved or almost always failed. Dropping the division restores plain
accuracy as the target. MaxRL (Tajwar et al., 2026) instead scales each group by its mean
success, deliberately emphasising hard prompts, and motivates this from a likelihood view of
success. Each choice can be worked out from the large-$G$ argument above. What matters for
engineering is that \emph{changing the normaliser is changing the objective}. Compare
variants at equal compute with several seeds, as \cref{sec:pg:loop} recommends, and never
treat the normaliser as a free implementation detail.

\begin{pitfall}[finite groups are not the limit]
The large-$G$ argument tells you what the estimator is aiming at, not what it does with
$G=8$. With small groups, $\bar r_i$ and the group's spread are noisy and correlated with
the very sample being weighted. Groups with zero spread also need an $\epsilon$ in any
denominator, or they produce NaNs. Implementations add $\epsilon$ to the standard deviation
(typically $10^{-6}$--$10^{-4}$); check what yours does when all $G$ rewards are equal (the
advantage must be exactly zero).
\end{pitfall}

\subsection{Per-sequence averaging reweights tokens}\label{sec:pg:length}

The second change is how per-token terms become a scalar. Three aggregations are in use
(\cref{fig:pg:weights}): average within each sequence and then over sequences
(\emph{sequence mean}, the original GRPO); average over all tokens in the batch
(\emph{token mean}, DAPO); and divide the sum by a fixed constant $Z$, such as batch size
times maximum length (\emph{constant}, Dr.\ GRPO). They give each token of response $i$ the
weights
\begin{equation}
\text{seq-mean: }\frac{1}{N\,|y_i|},\qquad
\text{token-mean: }\frac{1}{\sum_k|y_k|},\qquad
\text{constant: }\frac{1}{Z}.
\label{eq:pg:tokweights}
\end{equation}
Only the last two give every token the same weight; only the last is independent of the
batch. Under the sequence mean, a response's total push is independent of its length, so each
token of a long response is pushed less. If the advantage is negative, a long wrong answer is
penalised less per token than a short wrong answer, which can reward verbosity in failure;
\cref{ch:grpo2} traces the empirical consequences. The engineering consequence is that
aggregation must be implemented exactly as intended. \Cref{sec:pg:aggregation} shows that
gradient accumulation silently changes it if you are careless.

\begin{figure}[t]
\centering
\begin{tikzpicture}
\begin{axis}[ybar,width=0.85\linewidth,height=4.6cm,bar width=9pt,
  symbolic x coords={seq-mean,token-mean,constant},xtick=data,
  ylabel={weight per token ($\times10^{-3}$)},ymin=0,ymax=13,
  legend style={font=\scriptsize,at={(0.98,0.97)},anchor=north east},
  tick label style={font=\small},label style={font=\small}]
\addplot[fill=sysblue!60] coordinates {(seq-mean,12.5) (token-mean,1.10) (constant,0.98)};
\addplot[fill=napkin!70] coordinates {(seq-mean,0.31) (token-mean,1.10) (constant,0.98)};
\legend{response of 20 tokens, response of 800 tokens}
\end{axis}
\end{tikzpicture}
\caption{Weight that one token receives under the three aggregations of \cref{eq:pg:tokweights},
for a batch of $N=4$ responses with lengths $20,\,20,\,70,\,800$ (910 tokens) and constant
$Z=N\cdot256=1024$. Sequence-mean weights tokens of short responses $40\times$ more than
tokens of the long one; the other two weight all tokens equally.}
\label{fig:pg:weights}
\end{figure}

% =====================================================================
\section{From text to tensors: the RL batch}\label{sec:pg:batch}
\covers{CS336 A5 §4.2.1}

Everything so far assumed we can evaluate $\log\pitheta(y_t\mid x,y_{<t})$ for the sampled
tokens. Doing this correctly, and in bounded memory, is the first real engineering task of an
RL trainer.

\paragraph{Keep the sampled token IDs.} A rollout engine returns both text and token IDs.
\emph{Train on the IDs.} Re-tokenising the decoded text is a classic silent bug: BPE merges
across the prompt/response boundary (for example, a trailing space in the prompt merging with
the first word of the response), and normalisation of whitespace or Unicode can produce a
different token sequence from the one the policy actually sampled. The log-probabilities then
belong to a sequence that was never generated, and the gradient is wrong by an amount that
depends on the tokenizer. If you must work from strings, for example because the grader
rewrote the answer, tokenise prompt and response \emph{separately} and concatenate the ID
lists, so the boundary is explicit.

\paragraph{Layout.} For causal-LM training each row is \texttt{prompt + response}, right-padded.
The model reads positions $0..T-2$ and predicts positions $1..T-1$, so inputs and labels
are the same row shifted by one. A boolean \emph{response mask} aligned with the
\emph{labels}, not the inputs, marks which predictions are response tokens; the prompt and
the padding contribute nothing (\cref{fig:pg:layout}). Getting the alignment off by one is the
second classic silent bug. It trains the model on the last prompt token and drops the last
response token, and the loss still decreases.

\begin{figure}[t]
\centering
\begin{tikzpicture}[c/.style={draw,minimum width=7.6mm,minimum height=5.6mm,font=\scriptsize,inner sep=0pt}]
\def\row#1#2#3{%
 \foreach \v/\k [count=\i from 0] in {#3} {
   \ifnum\k=0 \def\col{fdata!25}\fi \ifnum\k=1 \def\col{fcompute!25}\fi \ifnum\k=2 \def\col{black!8}\fi
   \node[c,fill=\col] at (\i*0.8,#1) {\v};}
 \node[anchor=east,font=\scriptsize] at (-0.5,#1) {#2};}
\row{0}{full row}{p1/0,p2/0,p3/0,r1/1,r2/1,r3/1,pad/2}
\row{-0.9}{\code{input\_ids}}{p1/0,p2/0,p3/0,r1/1,r2/1,r3/1}
\row{-1.6}{\code{labels}}{p2/0,p3/0,r1/1,r2/1,r3/1,pad/2}
\foreach \v [count=\i from 0] in {0,0,1,1,1,0} {
  \ifnum\v=1 \node[c,fill=keygreen!35] at (\i*0.8,-2.4) {1};\else \node[c,fill=white] at (\i*0.8,-2.4) {0};\fi}
\node[anchor=east,font=\scriptsize] at (-0.5,-2.4) {\code{response\_mask}};
\draw[-{Stealth[length=1.5mm]},warnred] (1.6,-1.85)--(1.6,-2.15);
\node[font=\scriptsize,warnred,anchor=west,align=left] at (5.3,-2.0) {mask aligned with \emph{labels}:\\position 2 predicts r1, the first\\response token, so it is 1};
\node[font=\scriptsize,anchor=west] at (5.3,-0.2) {\fcolorbox{black}{fdata!25}{\phantom{x}} prompt\quad
 \fcolorbox{black}{fcompute!25}{\phantom{x}} response\quad \fcolorbox{black}{black!8}{\phantom{x}} pad};
\end{tikzpicture}
\caption{The RL training batch for one row. Inputs drop the final position and labels drop the
first. The response mask is defined on the label grid: it is 1 wherever the \emph{predicted}
token belongs to the response. Padding is ignored by the mask, so its ID never matters.}
\label{fig:pg:layout}
\end{figure}

\begin{lstlisting}[caption={Building the RL batch from token-ID lists (\code{llmsys.pg.build\_rl\_batch}).},label={lst:pg:batch}]
def build_rl_batch(prompts, responses, pad_id):
    full = [list(p) + list(r) for p, r in zip(prompts, responses)]
    T = max(len(s) for s in full)
    ids = torch.full((len(full), T), pad_id, dtype=torch.long)
    resp = torch.zeros((len(full), T), dtype=torch.bool)
    for i, (p, s) in enumerate(zip(prompts, full)):
        ids[i, :len(s)] = torch.tensor(s)
        resp[i, len(p):len(s)] = True               # positions holding response tokens
    return {"input_ids": ids[:, :-1],               # read 0..T-2
            "labels": ids[:, 1:],                   # predict 1..T-1
            "response_mask": resp[:, 1:]}           # aligned with labels
\end{lstlisting}

\subsection{Log-probabilities and entropy without a 60\,GB tensor}\label{sec:pg:logprobs}

Given logits $z\in\R^{N\times T\times V}$, the per-token log-probability of the label is
$z_{\text{label}}-\operatorname{logsumexp}(z)$, and the entropy of the next-token distribution
is
\begin{equation}
H=-\sum_v p_v\log p_v=-\sum_v p_v\big(z_v-\operatorname{logsumexp}(z)\big)
 =\operatorname{logsumexp}(z)-\sum_v p_v z_v ,
\label{eq:pg:entropy}
\end{equation}
which needs one softmax and no separate log-softmax tensor. Entropy is the single most useful
health metric of an RL run (\cref{ch:grpo2}): a collapse towards zero means the policy has
stopped exploring.

\begin{napkinmath}[the logits are the biggest tensor in RL]
A rollout batch of $N=256$ responses with $T\approx600$ tokens (prompt plus response) and a
vocabulary of $V=100{,}278$ (OLMo~2's) has
$256\times600\times100{,}278\approx1.54\times10^{10}$ logits: $30.8$\,GB in bf16 and
$61.6$\,GB if upcast to fp32 for a numerically safe log-softmax. That is more than the
weights, gradients and optimizer state of a 1B-parameter policy combined. Three remedies
compose. (i) \emph{Micro-batch}: 8 sequences at a time cuts the fp32 logits to $1.9$\,GB.
(ii) \emph{Chunk the vocabulary reduction}: process rows of the flattened $(N\cdot T)\times V$
matrix in slices, so only a slice is ever upcast; 2048 rows cost $0.82$\,GB. (iii) \emph{Never
materialise logits for prompt tokens}: slice the hidden states to response positions before
the LM head. A fused ``linear plus cross-entropy'' kernel goes further and never writes the
logits to HBM at all.
\end{napkinmath}

\begin{lstlisting}[caption={Chunked per-token log-probabilities and entropies in fp32 (\code{llmsys.pg.token\_logprobs\_entropy}).},label={lst:pg:logprobs}]
def token_logprobs_entropy(logits, labels, chunk_rows=None, with_entropy=True):
    B, T, V = logits.shape
    flat, lab = logits.reshape(B * T, V), labels.reshape(B * T)
    step = chunk_rows or flat.shape[0]
    lps, ents = [], []
    for s in range(0, flat.shape[0], step):
        z = flat[s:s + step].float()                  # upcast one slice only
        lse = torch.logsumexp(z, -1)
        lps.append(z.gather(-1, lab[s:s + step, None]).squeeze(-1) - lse)
        if with_entropy:
            ents.append(lse - (torch.softmax(z, -1) * z).sum(-1))   # Eq. for H
    lp = torch.cat(lps).view(B, T)
    return lp, (torch.cat(ents).view(B, T) if with_entropy else None)
\end{lstlisting}

\begin{pitfall}[log-probabilities of which policy?]
If the rollout engine samples at temperature $\tau\neq1$ or with top-$p$/top-$k$, the
\emph{behaviour} policy is the tempered or truncated distribution, not $\softmax(z)$. Old
log-probabilities used for importance ratios must then be computed from the same
transformed logits (divide by $\tau$, renormalise over the kept set), or every ratio is
biased before training even begins. Most LLM-RL recipes sample with $\tau=1$ and no
truncation for exactly this reason; \cref{ch:rlsys} treats the remaining engine/trainer
numerical mismatch.
\end{pitfall}

% =====================================================================
\section{Aggregation and gradient accumulation, done exactly}\label{sec:pg:aggregation}
\covers{CS336 A5 §4.2.3--4.2.4}

A rollout batch of $256$ long responses does not fit through the backward pass at once, so the
trainer splits it into micro-batches and accumulates gradients (\cref{ch:pretrain}). The
requirement is simple to state: the accumulated gradient must equal the gradient of the
full-batch loss. Whether it does depends on the aggregation.

\begin{proposition}[exact accumulation]\label{prop:pg:accum}
Let the full-batch loss be $\mathcal{L}=\frac{1}{D}\sum_{(i,t)\in\mathcal{M}}c_{it}\,\ell_{it}$
for per-token terms $\ell_{it}$, per-token coefficients $c_{it}$ that depend only on the row
$i$, and a normaliser $D$ computed over the \emph{whole} batch. If the batch is split into
micro-batches $\mathcal{M}_1,\dots,\mathcal{M}_K$ and each backward pass uses
$\mathcal{L}_k=\frac{1}{D}\sum_{(i,t)\in\mathcal{M}_k}c_{it}\ell_{it}$ with the \emph{same}
$D$, then $\sum_k\nabla\mathcal{L}_k=\nabla\mathcal{L}$ exactly, for any split.
\end{proposition}
The proposition is linearity, but it rules out the tempting implementations. For
\emph{sequence mean}, $c_{it}=1/|y_i|$ is row-local and $D=N$ is the total number of
sequences: divide by $N$, not by the micro-batch size. For \emph{token mean},
$c_{it}=1$ and $D$ is the total number of response tokens in the whole batch, which must
be counted \emph{before} the first backward pass. Dividing by each micro-batch's own token
count weights tokens differently in different micro-batches, and the error depends on how
the batch happened to be split. For \emph{constant}, $D=Z$ is fixed and accumulation is
trivially exact. \Cref{lst:pg:aggregate} takes the global normalisers as arguments; Lab~24
tests that one micro-batch and four give identical gradients for all three modes.

\begin{lstlisting}[caption={Aggregation that is exact under accumulation (\code{llmsys.pg.aggregate\_loss}).},label={lst:pg:aggregate}]
def aggregate_loss(per_token, mask, mode="seq_mean", n_seq_total=None,
                   n_tok_total=None, Z=None):
    m = mask.to(per_token.dtype)
    if mode == "seq_mean":       # mean over each row, then sum / (#rows in FULL batch)
        per_seq = (per_token * m).sum(1) / m.sum(1).clamp_min(1.0)
        return per_seq.sum() / (n_seq_total or per_token.shape[0])
    if mode == "token_mean":     # sum / (#response tokens in FULL batch)
        return (per_token * m).sum() / (n_tok_total or m.sum().clamp_min(1.0))
    if mode == "constant":       # sum / fixed Z
        return (per_token * m).sum() / Z
    raise ValueError(mode)
\end{lstlisting}

\begin{systemslens}[global normalisers under data parallelism]
With data parallelism (\cref{ch:distributed,ch:sharding}), ``the whole batch'' spans every rank.
DDP \emph{averages} gradients across $W$ ranks, which is correct only if each rank's loss was
normalised by the global count divided by $W$. If ranks hold different numbers of response
tokens, as they always do with variable-length rollouts, compute the global token count with
one all-reduce of a scalar \emph{before} the backward passes, divide by it, and multiply by
$W$ to undo DDP's averaging. The cost is one 4-byte collective per optimizer step; forgetting
it gives ranks with short responses too much influence.
\end{systemslens}

\paragraph{Skipping work that contributes nothing.} By \cref{prop:pg:accum}, a row whose
coefficients are all zero contributes nothing, so it can be dropped \emph{as long as} $D$
is computed from the original batch. With group baselines and binary rewards this is
common (see the napkin math of \cref{sec:pg:practical}), and with rejection fine-tuning
(\cref{sec:pg:rft}) every failed response qualifies. Pruning before the forward pass saves
both the forward and the backward cost of those rows. Re-pack the survivors into fewer
micro-batches to keep the GPU busy.

% =====================================================================
\section{Reusing samples: importance weights, TRPO and PPO}\label{sec:pg:offpolicy}
\covers{PO §3.3--3.4; PO §4.2--4.3; PO §4.5--4.6; CS336 A5 §6; CS336 L16}

Rollouts are the expensive part of LLM-RL: decoding is memory-bound and suffers from the long
tail (\cref{ch:inference,ch:rlsys}). Taking a single gradient step per rollout batch leaves
the trainer idle and the samples underused. Taking several steps per batch makes every step
after the first \emph{off-policy}: the data came from $\piold$, not from the current
$\pitheta$. This section derives what to do about it.

\subsection{Exact importance weighting, and why it fails for text}\label{sec:pg:is}

\Cref{eq:pg:is} corrects the expectation exactly. For a response the weight is a product over
tokens,
\begin{equation}
\frac{\pitheta(y\mid x)}{\piold(y\mid x)}=\prod_{t=1}^{|y|}w_t,\qquad
w_t=\frac{\pitheta(y_t\mid x,y_{<t})}{\piold(y_t\mid x,y_{<t})},
\end{equation}
and products of many random factors have enormous variance. Model the per-token ratios as
independent with $\E[w_t]=1$ and $\E[w_t^2]=1+\chi^2$, where $\chi^2=\sum_v p_1(v)^2/p_0(v)-1$ is
the chi-squared divergence between the new and old next-token distributions. Then
\begin{equation}
\operatorname{Var}\Big[\prod_{t=1}^{L}w_t\Big]=(1+\chi^2)^L-1 ,
\end{equation}
exponential in the response length. Even a tiny per-token divergence becomes catastrophic
at reasoning lengths (\cref{fig:pg:isvar}). With $\chi^2=0.01$, the variance is $0.10$ at
$L=10$, $1.7$ at $L=100$ and about $2\times10^4$ at $L=1000$. A handful of sequences with huge
weights then dominate every update.

\begin{figure}[t]
\centering
\begin{tikzpicture}
\begin{axis}[width=0.78\linewidth,height=5cm,xlabel={response length $L$ (tokens)},
  ylabel={$\log_{10}$ variance of the IS weight},xmin=0,xmax=1000,ymin=-2.5,ymax=4.8,grid=major,
  grid style={black!10},legend style={font=\scriptsize,at={(0.03,0.97)},anchor=north west},
  tick label style={font=\scriptsize},label style={font=\small}]
\addplot[thick,warnred] coordinates {(1,-2.000) (50,-0.191) (100,0.232) (200,0.800) (300,1.274)
 (400,1.720) (500,2.158) (600,2.592) (700,3.025) (800,3.457) (900,3.889) (1000,4.321)};
\addplot[thick,keygreen,dashed] coordinates {(1,-2) (1000,-2)};
\legend{sequence ratio $\prod_t w_t$, single token ratio $w_t$}
\end{axis}
\end{tikzpicture}
\caption{Variance of the exact sequence-level importance weight versus a single token's
weight, for i.i.d.\ tokens with $p_0=(0.5,0.3,0.2)$ and $p_1=(0.45,0.33,0.22)$
($\chi^2=0.01$). The sequence weight's variance grows as $(1+\chi^2)^L-1$. Computed with
\code{llmsys.pg.seq\_is\_variance} and checked by Monte Carlo in Lab~24.}
\label{fig:pg:isvar}
\end{figure}

\paragraph{The token-level surrogate.} PPO and GRPO replace the product by per-token ratios:
each token's log-probability gradient is weighted by its own $w_t$,
\[
\hat g_{\text{tok}}=\frac{1}{N}\sum_i A_i\sum_t w_{it}\,\nabla\log\pitheta(y_{it}\mid\cdot).
\]
This is biased: it reweights the choice of token $t$ but not the probability of the prefix
that led there, nor of the suffix that followed. As Degris et al.\ (2012) showed for
off-policy actor--critic methods, it optimises a surrogate in which the context of each
decision is still distributed as under the old policy. The bias is small when $\pitheta$
stays close to $\piold$. That observation motivates every trust-region method that follows,
and GSPO's geometric-mean sequence ratio (\cref{ch:grpo2}) is one way to restore
sequence-level reweighting without the exponential variance.

\subsection{TRPO: a trust region in KL}\label{sec:pg:trpo}

Trust Region Policy Optimization (Schulman et al., 2015) makes ``stay close'' explicit. It
maximises the importance-weighted surrogate subject to a bound on the expected KL divergence:
\begin{equation}
\max_\theta\;\E_{s,a\sim\piold}\Big[\frac{\pitheta(a\mid s)}{\piold(a\mid s)}A^{\piold}(s,a)\Big]
\quad\text{s.t.}\quad \E_s\big[\KL\big(\piold(\cdot\mid s)\,\|\,\pitheta(\cdot\mid s)\big)\big]\le\delta .
\end{equation}
Its theory bounds how much true performance can drop when the surrogate improves inside the
trust region. Its practice linearises the objective and quadratically approximates the
constraint with the Fisher information matrix $F$. The step is the \emph{natural gradient}
$\Delta\theta\propto F^{-1}g$, found with conjugate gradient using Fisher-vector products, then
scaled to the KL boundary and verified by a line search. For an LLM with $10^9$--$10^{11}$
parameters, each conjugate-gradient iteration costs a forward and backward pass, the line
search costs more forward passes, and the method composes poorly with Adam, mixed precision
and sharded optimizers. TRPO survives in LLM practice mainly through its ideas: ratios, trust
regions, and KL to the old policy.

\subsection{PPO: clipping instead of constraining}\label{sec:pg:ppo}

Proximal Policy Optimization (Schulman et al., 2017) replaces the constraint with a
first-order objective that simply stops rewarding movement beyond a band around
$w=1$:
\begin{equation}
\mathcal{J}^{\text{CLIP}}(\theta)=\E_t\Big[\min\big(w_t A_t,\;\clip(w_t,1-\epsilon,1+\epsilon)\,A_t\big)\Big],
\qquad \epsilon\approx0.1\text{--}0.3 .
\label{eq:pg:ppo}
\end{equation}
The $\min$ makes the objective a pessimistic bound. Differentiating it case by case
(\cref{fig:pg:clipmask}) gives the gradient that is actually applied:
\begin{equation}
\nabla\mathcal{J}^{\text{CLIP}}=\E_t\big[\,m_t\,w_t A_t\,\nabla\log\pitheta(a_t\mid s_t)\big],\qquad
m_t=\begin{cases}\mathbb{1}[w_t<1+\epsilon] & A_t\ge0\\ \mathbb{1}[w_t>1-\epsilon] & A_t<0.\end{cases}
\end{equation}
A token whose probability has already risen by more than a factor $1+\epsilon$ (good action)
or fallen below $1-\epsilon$ (bad action) stops receiving gradient. Movement in the
\emph{corrective} direction is never clipped. The fraction of masked tokens, the \emph{clip
fraction}, is the most direct gauge of how off-policy an update has become. A value near zero
means the extra epochs barely moved the policy; a large value means most of the batch has
stopped contributing.

\begin{figure}[t]
\centering
\begin{tikzpicture}[x=6.2cm,y=1.25cm]
\fill[keygreen!18] (0.3,0) rectangle (1.2,1.4);
\fill[warnred!18] (1.2,0) rectangle (1.75,1.4);
\fill[warnred!18] (0.3,-1.4) rectangle (0.8,0);
\fill[keygreen!18] (0.8,-1.4) rectangle (1.75,0);
\draw[->] (0.25,0)--(1.8,0) node[right,font=\small]{$w_t=\pitheta/\piold$};
\draw[->] (0.3,-1.5)--(0.3,1.55) node[above,font=\small]{$A_t$};
\draw[dashed] (1,-1.4)--(1,1.4) node[above,font=\scriptsize]{$1$};
\draw[dashed] (0.8,-1.4)--(0.8,1.4) node[above,font=\scriptsize]{$1-\epsilon$};
\draw[dashed] (1.2,-1.4)--(1.2,1.4) node[above,font=\scriptsize]{$1+\epsilon$};
\node[font=\scriptsize,align=center] at (0.53,0.7) {$A>0$, gradient flows:\\raise $\pi(a)$};
\node[font=\scriptsize,align=center] at (1.48,0.7) {already raised\\enough: $m_t=0$};
\node[font=\scriptsize,align=center] at (0.55,-0.7) {already lowered\\enough: $m_t=0$};
\node[font=\scriptsize,align=center] at (1.47,-0.7) {$A<0$, gradient flows:\\lower $\pi(a)$};
\end{tikzpicture}
\caption{Where PPO's clipped objective passes gradient (green) and where it masks it (red), by
importance ratio and sign of the advantage. Clipping only stops a token from moving
\emph{further} in the direction the advantage already pushed it.}
\label{fig:pg:clipmask}
\end{figure}

\paragraph{Why PPO is still an on-policy method.} The data are generated by the current policy
and thrown away after $K$ epochs (typically $K=1$--$4$); clipping only makes those few
reuses safe. PPO cannot learn from an arbitrary replay buffer or an old dataset. That
distinguishes it from genuinely off-policy methods, and from offline preference methods such
as DPO, which the comparison in \cref{tab:pg:compare} places side by side.

\paragraph{The rest of the PPO recipe.} Three extra terms are common. An \emph{entropy bonus}
$+c_H\,H(\pitheta(\cdot\mid s_t))$ discourages premature collapse; it is mostly omitted for
LLMs, where the vocabulary already provides plenty of entropy. A \emph{KL penalty} to a
reference policy, either in the loss ($-\beta\,\KL(\pitheta\|\piref)$, GRPO's form) or folded
into a per-token reward (PPO-RLHF's form), anchors the policy to its starting point;
\cref{ch:grpo2} implements the estimators. PPO's original paper also proposed an adaptive
KL-\emph{penalty} variant with no clipping, which adjusts $\beta$ to hit a KL target. A
\emph{value loss} trains the critic; \cref{ch:rlhf} develops it together with generalised
advantage estimation.

\begin{table}[ht]
\centering\small
\begin{tabularx}{\linewidth}{@{}lXXXX@{}}
\toprule
& REINFORCE / RLOO & TRPO & PPO & GRPO\\
\midrule
Data & online, on-policy & online, on-policy & online, $K$ reuses & online, $K$ reuses\\
Baseline & none / group LOO & critic $V$ & critic $V$ + GAE & group mean (and std)\\
Update control & step size & KL constraint & ratio clipping & ratio clipping + KL loss\\
Extra networks & none & critic & critic & reference (optional)\\
Per-step cost & 1 fwd+bwd & CG + line search & $K\times$ fwd+bwd & $K\times$ fwd+bwd\\
Variance & high & low & low & medium\\
LLM use & RLOO, RFT & rare & RLHF classic & RLVR standard\\
\bottomrule
\end{tabularx}
\caption{The policy-gradient family at a glance. Preference-based offline methods (DPO, KTO)
are compared with these in \cref{ch:dpo}.}
\label{tab:pg:compare}
\end{table}
% =====================================================================
\section{Learning only from successes: rejection fine-tuning}\label{sec:pg:rft}
\covers{CS336 A5 §5.2; CS336 L16}

The simplest way to improve a model with a verifier needs no policy-gradient machinery at all.
Sample several responses per prompt, keep the ones the verifier accepts, and fine-tune on them
with the ordinary next-token loss. Repeat with the improved model. This family goes by many
names: \emph{rejection-sampling fine-tuning} (RFT), \emph{expert iteration}, STaR (Zelikman et
al., 2022, which adds rationalisation hints for failed problems), ReST (Gulcehre et al., 2023,
which separates a ``grow'' phase of sampling from an ``improve'' phase of repeated filtering and
training), and RAFT (Dong et al., 2023). Its gradient for one round is
\begin{equation}
\hat g_{\text{RFT}}=\frac{1}{Z}\sum_{i,j}\mathbb{1}\big[r(x_i,y_{ij})=1\big]\,\nabla_\theta\log\pitheta(y_{ij}\mid x_i),
\end{equation}
which, for 0/1 rewards, is exactly the REINFORCE estimator \cref{eq:pg:reinforce} with
\emph{no baseline} and a constant normaliser. On our toy task (\cref{fig:pg:curves}) the two
produce the same curve to the last digit when trained with Adam, because they produce the same
gradient up to a constant factor.

RFT's appeal is mostly systems simplicity. A round can be fully offline: generate a large
pool with an inference engine, filter it with a verifier, and hand the survivors to an
ordinary SFT trainer, with no weight synchronisation, no old log-probabilities and no
importance weights. Failures never enter the trainer at all, so for hard tasks, where most
samples fail, the training cost per round is tiny. What it gives up is the negative signal.
A baseline-subtracted estimator also pushes \emph{down} the probability of failed responses,
which RFT never does, so RFT tends to reduce diversity faster: it can only concentrate mass on
what already works. In practice RFT is used as a warm start, as a distillation-style data
generator (\cref{ch:distill}), or as a strong, cheap baseline against which an RL method must
justify its extra machinery.

% =====================================================================
\section{The engineering loop}\label{sec:pg:loop}
\covers{CS336 A5 §4.2.2, §4.3; CS336 L16}

We now have every mathematical piece. This section assembles them into the system that
\cref{ch:grpo,ch:grpo2} use on a single GPU, that \cref{ch:rlsys} scales out, and that the
CS336 assignment (\cref{ch:proj5}) asks you to build with a 1B-parameter model on two GPUs.

\subsection{Architecture: who runs where}

An RL step has four stages: \emph{generate} rollouts, \emph{score} them, \emph{compute
advantages}, and \emph{update} the policy (\cref{fig:pg:system}). The first is inference and
the last is training, and they want different software. Inference wants an engine with paged
KV caches, continuous batching and CUDA graphs (\cref{ch:serving}). Training wants autograd,
an optimizer and gradient accumulation. The simplest real deployment puts them on different
GPUs.
\begin{itemize}
\item \textbf{GPU 0, the trainer}: the policy in bf16 with fp32 optimizer state, computing
  log-probabilities and the loss and taking optimizer steps.
\item \textbf{GPU 1, the rollout engine}: an inference engine (vLLM, SGLang, or the engine of
  \cref{ch:serving}) holding its own copy of the weights plus a large KV-cache pool.
\item \textbf{Weight sync}: after each optimizer step, the trainer's weights are broadcast to
  the engine over NCCL (a dedicated process group between the two), and the engine
  invalidates caches that depended on the old weights, such as the prefix cache.
\item \textbf{Reward}: a CPU function (regex and SymPy, \cref{ch:eval}) or a reward-model
  server; sandboxed if it executes code (\cref{ch:rlsys}).
\end{itemize}
The alternative \emph{colocated} layout time-shares one GPU: the engine releases or offloads
its KV cache while the trainer runs. \Cref{ch:rlsys} gives the cost model for choosing between
the two.

\begin{figure}[t]
\centering
\resizebox{\linewidth}{!}{%
\begin{tikzpicture}[b/.style={draw,rounded corners=2pt,align=center,font=\small,minimum height=9mm,text width=26mm},
  >={Stealth[length=2mm]}]
\node[b,fill=fdata!15] (data) at (0,0) {prompt sampler\\$B$ prompts};
\node[b,fill=fmemory!25,text width=30mm] (eng) at (4.4,0) {rollout engine\\{\scriptsize GPU 1: weights + KV pool}};
\node[b,fill=napkin!20] (rew) at (8.8,0) {reward fn\\{\scriptsize CPU / sandbox / RM}};
\node[b,fill=fcomm!15] (adv) at (13.0,0) {advantages\\{\scriptsize group statistics}};
\node[b,fill=fcompute!20,text width=34mm] (tr) at (8.8,-2.6) {trainer {\scriptsize(GPU 0)}\\log-probs $\to$ loss $\to$\\accumulate $\to$ clip $\to$ step};
\node[b,fill=black!5,text width=30mm] (log) at (13.0,-2.6) {logger\\{\scriptsize reward, entropy, grad norm,\\length, clip fraction}};
\draw[->,thick] (data)--node[above,font=\scriptsize]{$x\times G$}(eng);
\draw[->,thick] (eng)--node[above,font=\scriptsize]{token IDs}(rew);
\draw[->,thick] (rew)--node[above,font=\scriptsize]{$r$}(adv);
\draw[->,thick] (adv.south) -- ++(0,-0.7) -| ([xshift=6mm]tr.north);
\draw[->,thick] (eng.south) |- node[pos=0.3,left,font=\scriptsize,align=right]{token IDs\\(+ old log-probs)} (tr.west);
\draw[->,thick,dashed] (tr.east)--(log.west);
\draw[->,very thick,warnred] (tr.south) -- ++(0,-0.7) -| node[pos=0.25,below,font=\scriptsize,warnred]{NCCL broadcast of new weights after every optimizer step} (eng.south);
\end{tikzpicture}}
\caption{The minimal two-GPU RL system. Rollouts flow left to right as token IDs, never as
re-tokenised text; the trainer pushes new weights back to the engine after each optimizer
step. With $K>1$ updates per batch, the old log-probabilities come either from the engine
(cheap, but numerically different) or from a trainer forward pass before the first update
(exact, one extra forward).}
\label{fig:pg:system}
\end{figure}

\subsection{One step, end to end}

\Cref{lst:pg:step} is the whole loop for the toy task of Lab~24, written so that every line
has a counterpart at scale: \code{sample\_fixed\_length} becomes a call to the engine,
\code{task.reward} becomes the verifier, and the inner loop becomes micro-batched
accumulation.

\begin{lstlisting}[caption={One on-policy step with optional off-policy reuse (abridged from \code{llmsys.pg.train\_toy}).},label={lst:pg:step}]
prompts = task.sample_prompts(B, gen).repeat_interleave(G, 0)    # (B*G, P)
resp, old_lp = sample_fixed_length(model, prompts, task.k, generator=gen)
r = task.reward(prompts, resp)                                    # 0/1 per row
adv = group_advantages(r.view(B, G), "mean", "std").reshape(-1)
method = "ppo_clip" if n_minibatches > 1 else "reinforce"         # reuse => clip
for sel in torch.randperm(B * G).chunk(n_minibatches):
    lp, ent = response_logprobs(model, prompts[sel], resp[sel])   # current policy
    mask = torch.ones_like(lp, dtype=torch.bool)
    per_tok, stats = pg_token_loss(lp, adv[sel], mask, old_lp[sel], method)
    loss = aggregate_loss(per_tok, mask, "seq_mean")
    opt.zero_grad(set_to_none=True)
    loss.backward()
    grad_norm = torch.nn.utils.clip_grad_norm_(model.parameters(), 1.0)
    opt.step()
# at scale: engine.sync_weights(model) here, before the next rollout
\end{lstlisting}

\Cref{fig:pg:curves} runs the chapter's estimators on the toy task (a 4-token answer whose
digit sum must hit a target residue; reward 1 or 0). Three lessons carry over to real
models. First, \emph{seeds matter}: one of three REINFORCE seeds without a baseline locked
into a 33\% plateau, while all baseline-subtracted seeds converged. Report several seeds,
never one. Second, the group baseline and leave-one-out runs are close, as the rescaling
identity of \cref{sec:pg:practical} predicts. Third, reusing each rollout batch for 8
minibatch updates reaches high reward in far fewer \emph{rollout batches}. On this easy
toy, the naive and clipped reuse runs both stayed stable; the instabilities that clipping
guards against appear on harder tasks and larger steps (\cref{ch:grpo2}). Reuse cost about
$2.8\times$ more trainer time per batch, a good trade whenever rollouts dominate the step
time, as they do for LLMs.

\begin{figure}[t]
\centering
\begin{tikzpicture}
\begin{axis}[width=0.9\linewidth,height=6cm,xlabel={rollout batch (16 prompts $\times$ 8 samples)},
  ylabel={training reward (mean of 3 seeds)},xmin=0,xmax=80,ymin=0,ymax=1.05,grid=major,
  grid style={black!10},legend style={font=\scriptsize,at={(0.98,0.03)},anchor=south east},
  tick label style={font=\scriptsize},label style={font=\small}]
\addplot[thick,napkin] coordinates {(2.5,0.104) (7.5,0.162) (12.5,0.236) (17.5,0.284) (22.5,0.321) (27.5,0.336) (32.5,0.352) (37.5,0.375) (42.5,0.426) (47.5,0.542) (52.5,0.605) (57.5,0.670) (62.5,0.697) (67.5,0.712) (72.5,0.755) (77.5,0.763)};
\addplot[thick,labpurple,dashed] coordinates {(2.5,0.107) (7.5,0.168) (12.5,0.231) (17.5,0.297) (22.5,0.306) (27.5,0.326) (32.5,0.382) (37.5,0.415) (42.5,0.470) (47.5,0.601) (52.5,0.701) (57.5,0.784) (62.5,0.855) (67.5,0.911) (72.5,0.958) (77.5,0.978)};
\addplot[thick,sysblue] coordinates {(2.5,0.108) (7.5,0.170) (12.5,0.236) (17.5,0.293) (22.5,0.317) (27.5,0.358) (32.5,0.401) (37.5,0.473) (42.5,0.548) (47.5,0.694) (52.5,0.795) (57.5,0.871) (62.5,0.948) (67.5,0.964) (72.5,0.975) (77.5,0.988)};
\addplot[thick,keygreen] coordinates {(2.5,0.184) (7.5,0.394) (12.5,0.530) (17.5,0.795) (22.5,0.891) (27.5,0.956) (32.5,0.994) (37.5,0.997) (42.5,0.998) (47.5,0.998) (52.5,0.999) (57.5,1.000) (62.5,0.998) (67.5,1.000) (72.5,0.998) (77.5,0.999)};
\addplot[thick,warnred,densely dotted] coordinates {(2.5,0.161) (7.5,0.309) (12.5,0.415) (17.5,0.543) (22.5,0.584) (27.5,0.713) (32.5,0.861) (37.5,0.963) (42.5,0.969) (47.5,0.808) (52.5,0.839) (57.5,0.910) (62.5,0.960) (67.5,0.994) (72.5,0.996) (77.5,0.996)};
\legend{{REINFORCE, no baseline (= RFT)},{leave-one-out},{group mean/std (GRPO)},{$8\times$ reuse, naive},{$8\times$ reuse, PPO-clip}}
\end{axis}
\end{tikzpicture}
\caption{Policy-gradient estimators on the digit-sum toy task (20k-parameter Qwen3-shaped
policy trained from scratch, Adam, learning rate $3\times10^{-3}$, 3 seeds, 5-batch moving
average). Measured on CPU in this environment with one thread, about 10\,s per on-policy run
and 28\,s per $8\times$-reuse run; reproduce with \code{run.py compare}. The REINFORCE mean
is pulled down by one of its three seeds, which stalled at 0.33.}
\label{fig:pg:curves}
\end{figure}

\subsection{Configuration, logging, and the run directory}

Treat every run as an experiment you will need to explain later. A minimal configuration
records everything that changes the estimator: the baseline, the advantage scaling, the
aggregation and its constant, the number of updates per rollout batch, the clip range,
the KL coefficient, the learning rate, the group size, the sampling temperature and maximum
length, and the seed. Pass the configuration as a dataclass or YAML file, never as code edits,
and write it to the run directory before the first step.

\begin{lstlisting}[style=plain,caption={A run directory that makes RL experiments reproducible and comparable.},label={lst:pg:rundir}]
runs/grpo-std-seqmean-lr1e-5/seed=2/
  config.json            # every knob + git commit + library versions
  metrics.jsonl          # one line per step (see the metrics table)
  rollouts/step_0040.jsonl   # prompts, responses, rewards: read them!
  ckpt/step_0200/        # weights + optimizer + RNG + data position
\end{lstlisting}

\begin{table}[ht]
\centering\small
\begin{tabularx}{\linewidth}{@{}lX@{}}
\toprule
Metric & What it tells you\\
\midrule
train / validation reward & the objective; validation on a fixed held-out set of at least $\sim10^3$ prompts\\
format reward & whether the model still emits parseable answers (separates format from correctness)\\
token entropy & exploration; a collapse towards 0 precedes stagnation\\
gradient norm (pre-clip) & spikes flag bad batches or numerical trouble\\
response length (mean, max, \% truncated) & length drift; truncation silently zeroes rewards\\
zero-advantage fraction & how much of the rollout compute is wasted (\cref{sec:pg:practical})\\
clip fraction, approx.\ KL to $\piold$ & how off-policy the reuse has become\\
pass@$k$ within a group & whether the group contains \emph{any} success to learn from\\
step-time breakdown & generate / score / train / sync, the input to every systems decision\\
\bottomrule
\end{tabularx}
\caption{Metrics every LLM-RL run should log. The first three are about the model; the rest are
about whether the estimator and the system behave as intended.}
\label{tab:pg:metrics}
\end{table}

\paragraph{Variance across runs.} RL runs from different seeds diverge far more than
supervised runs, because the data depend on the model. Plot every seed, or the mean with an
interval such as $\bar m\pm1.96\,s/\sqrt{n}$ over $n$ seeds. With $n=4$ that interval is
optimistic, so show the individual runs as well. Decide in advance how many seeds an ablation
needs, and change one thing at a time. Tuning the learning rate for the baseline method only
and then comparing variants at that rate establishes at best a lower bound for the variants.
\Cref{ch:proj5} turns this into a concrete plan. The generic helpers in
\code{code/projects/a5} expand a seed$\times$setting grid into commands and aggregate
per-seed logs into intervals.

\begin{napkinmath}[where does a step's time go?]
Take a 1.2B-parameter policy on H100s (\code{llmsys.flops.HARDWARE}), 256 rollouts per step,
prompts and responses averaging $\sim$400 tokens. \emph{Training}: $6ND=6\times1.2\times10^9
\times256\times400\approx7.4\times10^{14}$ FLOPs, or $1.9$\,s at 40\% MFU. \emph{Rollout}: one
decode step must at least read the 2.4\,GB of bf16 weights, $0.72$\,ms at 3.35\,TB/s, so 512
decode steps take at least $0.37$\,s even with all 256 sequences batched together. Real engines
spend several times that in launch overheads, sampling and KV traffic, and the slowest sequence
in the batch sets the pace. Rollout is therefore comparable to or larger than training even for
a small model, and the ratio worsens as responses grow. That is why reusing samples
(\cref{sec:pg:offpolicy}), pruning zero-advantage rows (\cref{sec:pg:aggregation}), and
asynchronous pipelines (\cref{ch:rlsys}) are worth their complexity.
\end{napkinmath}

\subsection{A debugging checklist}

When a run does not learn, check these in order. Each takes minutes and has caught real
bugs.
\begin{enumerate}
\item \textbf{Rewards}: print ten rollouts with their rewards and verify each by hand. A
  grader that never fires (wrong stop string, answer extraction off by one tag) looks
  exactly like a model that cannot learn.
\item \textbf{Masks}: assert \code{response\_mask.sum(1)} equals the response lengths, and that
  the first masked label equals the first response token (\cref{fig:pg:layout}).
\item \textbf{Log-probabilities}: on the first update, with no dropout, the trainer's
  log-probabilities must match the engine's sampling log-probabilities to within numerical
  tolerance. A large gap means a tokenisation, template or temperature mismatch.
\item \textbf{Advantages}: groups with identical rewards must give exactly zero advantage; the
  advantages of a group must sum to zero (mean baseline).
\item \textbf{Accumulation}: one micro-batch versus $k$ must give the same gradient
  (\cref{prop:pg:accum}); test it once with a tiny model.
\item \textbf{Weight sync}: after a sync, the engine's log-probability of a fixed string must
  change; if it does not, the engine is still sampling from the initial weights.
\item \textbf{Entropy and length}: a sudden entropy collapse or length explosion within a few
  steps usually means the learning rate is too high or the KL/clip settings were bypassed.
\end{enumerate}

% =====================================================================
\section{Summary}\label{sec:pg:summary}
\begin{itemize}
\item LLM generation is an MDP with deterministic transitions, a vocabulary-sized action space
  and (usually) a single terminal reward; \cref{tab:pg:notation} maps the vocabularies.
\item The score-function identity turns $\nabla\E[r]$ into $\E[r\nabla\log\pi]$; implementations
  minimise a surrogate whose value is meaningless and whose gradient is the estimator.
\item Baselines independent of the sample preserve the expectation and can cut variance by
  large factors; the optimal constant is a score-weighted mean, and leave-one-out is the
  group-mean baseline up to a factor $G/(G-1)$.
\item Scaling advantages by group statistics reweights prompts by difficulty, and per-sequence
  averaging reweights tokens by length: both change the objective.
\item Systems correctness: train on sampled token IDs, align the mask with labels, compute
  log-probs and entropy in chunked fp32, and normalise every micro-batch by global counts.
\item Sequence importance weights explode exponentially with length; PPO's clipped token ratios
  trade a controlled bias for bounded variance and give a clip fraction to monitor.
\item RFT is REINFORCE without a baseline on 0/1 rewards: cheap, offline-friendly, positive-only.
\item The working system is engine + reward + trainer + weight sync, instrumented with the
  metrics of \cref{tab:pg:metrics} and run over several seeds.
\end{itemize}

% =====================================================================
\section{Exercises}\label{sec:pg:exercises}

\begin{exercise}[optimal baseline, vector form]
Derive \cref{eq:pg:bstar} for a vector-valued score $g$ by minimising
$\E[(r-b)^2\|g\|^2]$ over $b$. Then generalise to a \emph{per-coordinate} baseline $b_k$ and
explain why it is never used for LLMs.
\end{exercise}

\begin{exercise}[the three-armed bandit by hand]
For $\pi=\softmax(0.5,0,-0.5)$ and $r=(1,0.2,0)$, compute the exact gradient with respect to the
logits, the total variance at $b=0$, and $b^\star$. Check your numbers against
\cref{fig:pg:variance} and \code{run.py bandit}.
\end{exercise}

\begin{exercise}[leave-one-out]
Prove that the leave-one-out baseline gives an unbiased estimator for any $G\ge2$, and show that
$r_j-\frac{1}{G-1}\sum_{k\ne j}r_k=\frac{G}{G-1}(r_j-\bar r)$.
\end{exercise}

\begin{exercise}[a new reweighting]
Using the large-$G$ argument of \cref{sec:pg:reweight}, find the prompt weight $w(x)$ induced by
the advantage $A_{ij}=(r_{ij}-\bar r_i)\cdot\bar r_i$ (scale by the group's success rate). Which
prompts does it emphasise? Sketch $w$ as a function of $\eta$.
\end{exercise}

\begin{exercise}[token weights]
A batch holds four responses of lengths 30, 30, 300 and 1500 with advantages $+1,+1,-1,-1$.
Compute the total gradient weight each response receives (sum over its tokens) under the three
aggregations of \cref{eq:pg:tokweights}, with $Z=4\cdot2048$. Which aggregation makes the long
wrong answer cheapest to keep?
\end{exercise}

\begin{exercise}[accumulation bugs]
A trainer uses token-mean aggregation but divides each micro-batch's summed loss by that
micro-batch's own token count, then averages over micro-batches. With two micro-batches
holding 100 and 900 response tokens, by what factor is each token's weight wrong in each
micro-batch, relative to the exact rule?
\end{exercise}

\begin{exercise}[importance-weight explosion]
For $p_0=(0.6,0.3,0.1)$ and $p_1=(0.55,0.33,0.12)$, compute $\chi^2$ and the variance of the
sequence importance weight at $L=256$ and $L=4096$. What would the effective sample size
$\frac{(\sum w)^2}{\sum w^2}$ of 1024 such sequences roughly be at $L=4096$?
\end{exercise}

\begin{exercise}[clip masks]
For $\epsilon=0.2$, state whether PPO passes gradient for each token:
$(w,A)=(1.3,+1),\,(1.3,-1),\,(0.7,+1),\,(0.7,-1),\,(1.1,+1),\,(0.9,-1)$. Explain the asymmetry.
\end{exercise}

\begin{exercise}[logit memory]
A 7B policy with $V=152{,}064$ trains on micro-batches of 4 sequences of 8192 tokens. How many
GB do fp32 logits for one micro-batch take? Choose a chunk size so that the fp32 working set
stays under 1\,GB, and estimate the cost of slicing to response positions only if responses are
75\% of tokens.
\end{exercise}

\begin{exercise}[measurement: estimator variance on the toy]
Use Lab~24's toy task to measure the empirical variance of the policy gradient (flattened over
parameters) for no baseline, group mean and leave-one-out, at initialisation and after 20
GRPO steps. Relate the change to the success rate $\eta$.
\end{exercise}

\begin{exercise}[design: zero-advantage pruning]
With $G=16$ and a dataset whose per-prompt success rates are uniform on $[0,1]$, what fraction
of groups has zero spread under 0/1 rewards? How much trainer time does pruning save, and why does
it save nothing on rollout time?
\end{exercise}

% =====================================================================
\section{Lab 24: Policy-gradient building blocks}\label{sec:pg:lab}
\begin{labbox}{24}{Policy-gradient building blocks}
\labdir{lab24_pg}
\labsection{Goals}
Implement, test and measure the estimator components of this chapter on problems small enough
to verify exactly, then watch them train a tiny policy from scratch.
\labsection{Setup}
\begin{lstlisting}[style=shell]
cd code
LAB_IMPL=starter python -m pytest -q labs/lab24_pg     # grade your work
python labs/lab24_pg/run.py bandit           # exact vs sampled
\end{lstlisting}
Pin PyTorch to one thread (\code{torch.set\_num\_threads(1)}) on a shared CPU: for a 20k-parameter
model it is many times faster than oversubscribing cores.
\labsection{Tasks}
\begin{enumerate}
\item \code{bandit\_estimates}, \code{optimal\_baseline}: sampled score-function estimates with no,
  constant and leave-one-out baselines; the variance-optimal constant.
\item \code{build\_rl\_batch}: the layout of \cref{fig:pg:layout}.
\item \code{token\_logprobs\_entropy}: chunked fp32 log-probs and entropy (\cref{lst:pg:logprobs}).
\item \code{group\_advantages}: mean / leave-one-out baselines, optional std scaling, safe on
  zero-spread groups.
\item \code{aggregate\_loss}: sequence-mean, token-mean, constant, with global normalisers.
\item \code{pg\_token\_loss}, \code{rft\_token\_loss}: REINFORCE, token IS, sequence IS, PPO-clip
  with clip fraction; rejection fine-tuning.
\item \code{accumulated\_gradient}: micro-batched REINFORCE gradient equal to the full-batch one
  for all three aggregations, with variable-length (padded) responses.
\item \code{seq\_is\_variance}: the closed form, checked by Monte Carlo.
\end{enumerate}
\labsection{Measurements}
\begin{itemize}
\item \code{run.py bandit}: reproduce \cref{fig:pg:variance}; confirm $b^\star<\E[r]$ here and explain why.
\item \code{run.py compare --seeds 3} then \code{run.py plot}: reproduce \cref{fig:pg:curves}. Record
  wall-clock per run and the reward per rollout batch versus per trainer second.
\item \code{run.py isvar}: tabulate the sequence-ratio variance against length.
\item Ablate the group size $G\in\{2,4,8,16\}$ at a fixed number of rollouts per step and plot final
  reward and zero-advantage fraction.
\end{itemize}
\labsection{Deliverables}
Passing tests; the three plots; one paragraph per plot explaining it with the chapter's
equations; the measured trainer-time cost of $8\times$ reuse.
\labsection{Stretch goals}
Add a learned scalar value head as a baseline (preview of \cref{ch:rlhf}) and compare its
gradient variance with leave-one-out; implement GSPO's geometric-mean sequence ratio and compare
its clip fraction with PPO's at matched $\epsilon$; port the trainer loop to Qwen3-0.6B with the
serving engine of \cref{ch:serving} as the rollout generator.
\end{labbox}
